\documentclass[11pt,a4paper]{article}
\usepackage[T1]{fontenc}
\usepackage[utf8]{inputenc}
\usepackage{lmodern,amsmath,amssymb}
\usepackage[margin=25mm]{geometry}
\usepackage{longtable,booktabs,array,calc}
\usepackage[hidelinks]{hyperref}
\hypersetup{pdftitle={The finiteness half of the mass gap is the nonvanishing of one smeared susceptibility},pdfauthor={navstokgap research working notes}}
\newcommand{\tightlist}{\setlength{\itemsep}{0pt}\setlength{\parskip}{0pt}}
\setlength{\emergencystretch}{3em}
\setcounter{secnumdepth}{0}
\begin{document}
\hypertarget{the-finiteness-half-of-the-mass-gap-is-the-nonvanishing-of-one-smeared-susceptibility}{%
\section{The finiteness half of the mass gap is the nonvanishing of one
smeared
susceptibility}\label{the-finiteness-half-of-the-mass-gap-is-the-nonvanishing-of-one-smeared-susceptibility}}

The requirement \(m<\infty\) of the Jaffe--Witten statement reduces to
an exact inequality between two ground-state quantities. For any
gauge-invariant local scalar \(\varphi\) built from the spatially
smeared field, with \(\langle\varphi\rangle=0\),
\[m\ \le\ \frac{g^2\hbar c}{2}\;\frac{\big\langle\,|D_\varphi|^2\,\big\rangle_0}{\chi_\varphi},
\qquad
D_\varphi=\int d^3y\;\frac{\delta\varphi(y)}{\delta A^a_i(0)},\qquad
\chi_\varphi=\int d^3z\;\big\langle\varphi(z)\varphi(0)\big\rangle_0^{\rm c},\]
in the continuum theory with a unique vacuum, whenever the finite-volume
gaps and correlators converge. Both sides are intensive: the volume
cancels between the numerator and the denominator, which is what makes
the zero-momentum operator the right trial excitation and the plaquette
of \href{lattice-gap-upper-bounds.md}{the upper-bound note} the wrong
one. The inequality is the energy-weighted sum rule, \(m\) bounded by
the mean excitation energy carried by \(\varphi\), so its content is:
\textbf{if some local gauge-invariant observable has nonvanishing
zero-momentum fluctuations in the vacuum, the theory has an excitation
of finite energy.} Conversely \(m=\infty\) forces every such
susceptibility to vanish, that is, the vacuum is an eigenstate of every
smeared gauge-invariant operator. The ultraviolet problem that defeated
the plaquette is solved by smearing, and the natural smearing is the
gradient flow, whose observables are renormalized; what is assumed, and
stated as hypothesis (F1), is that the equal-time flowed correlators
have continuum limits, which is not Lüscher's four-dimensional theorem
but its Hamiltonian counterpart. Dimensional analysis gives the
perturbative behaviour
\(\langle|D|^2\rangle/\chi\simeq c_\varphi/\sqrt{8t}\) at flow time
\(t\), so the bound reads \(m\le C g^2\hbar c/\sqrt{8t}\) while
perturbation theory applies; since the bound cannot fall below \(m\),
the flow time at which that scaling must break down is the correlation
length \(\hbar c/m\). In free Maxwell theory the scaling persists for
all \(t\) and the bound correctly tends to \(m=0\). The inequality
itself is the textbook single-mode (Feynman--Bijl) bound; what is
specific here is the zero-momentum intensive form, the flow as the
smearing that makes it cutoff-independent, and the identification with
the finiteness half of the Millennium statement. Constants explicit;
nothing promoted.

\hypertarget{the-continuum-bound}{%
\subsection{1. The continuum bound}\label{the-continuum-bound}}

Take the temporal-gauge Hamiltonian in the convention of
\href{low-dimensional-mass-gap.md}{G07} Proposition 7, with \(g\)
dimensionless and \([A]=L^{-1}\):
\[H=\int d^3x\;\Big[\frac{g^2c}{2\hbar}\,\Pi^a_i\Pi^a_i
+\frac{\hbar c}{4g^2}\,F^a_{ij}F^a_{ij}\Big],
\qquad\big[A^a_i(x),\Pi^b_j(y)\big]=i\hbar\,\delta^{ab}\delta_{ij}\delta^3(x-y),\]
obtained from
\(S=-\frac{\hbar}{4g^2}\int F^a_{\mu\nu}F^{a\mu\nu}d^3x\,c\,dt\) with
\(\Pi^a_i=(\hbar/g^2c)\dot A^a_i\).

\textbf{Proposition 1.} Let \(O\) be a real functional of the spatial
field \(A\) alone, gauge invariant, with \(\langle O\rangle_0=0\) and
\(O\Omega\in D(H)\). Then \[E_1-E_0\ \le\ \frac{g^2\hbar c}{2}\;
\frac{\displaystyle\int d^3x\;\Big\langle\Big|\frac{\delta O}{\delta A^a_i(x)}\Big|^2\Big\rangle_0}
{\big\langle O^2\big\rangle_0}.\]

\emph{Proof.} \(O\Omega\perp\Omega\), so by min-max
\(E_1-E_0\le\langle O\Omega,(H-E_0)O\Omega\rangle/\|O\Omega\|^2\), and
\(2\langle\Omega,O(H-E_0)O\Omega\rangle=\langle\Omega,[O,[H,O]]\Omega\rangle\)
since \([O,[H,O]]=2OHO-O^2H-HO^2\). The magnetic term commutes with
\(O\). For the electric term,
\([\Pi^a_i(x),O]=-i\hbar\,\delta O/\delta A^a_i(x)\) and
\[\big[O,[\Pi^a_i\Pi^a_i,O]\big]=-2\big[\Pi^a_i,O\big]^2
=2\hbar^2\Big|\frac{\delta O}{\delta A^a_i}\Big|^2,\] so
\([O,[H,O]]=\frac{g^2c}{2\hbar}\cdot2\hbar^2\int d^3x|\delta O/\delta A|^2 =g^2\hbar c\int d^3x|\delta O/\delta A|^2\).
Divide by \(2\langle O^2\rangle\). \(\square\)

Only the electric term contributes: the bound measures how much the
trial observable is disturbed by the conjugate field, which is the same
quantity that Section 2 of \href{action-floor-yang-mills-gap.md}{the
action-floor note} identifies as the carrier of \(\hbar\).

\hypertarget{zero-momentum-makes-the-bound-intensive}{%
\subsection{2. Zero momentum makes the bound
intensive}\label{zero-momentum-makes-the-bound-intensive}}

Let \(\varphi(x)\) be a gauge-invariant local scalar with
\(\langle\varphi\rangle_0=0\) and take \(O=\int_Vd^3y\,\varphi(y)\) in a
periodic box of volume \(V\). Write
\[D^a_i(x)=\frac{\delta O}{\delta A^a_i(x)}=\int_Vd^3y\;\frac{\delta\varphi(y)}{\delta A^a_i(x)} .\]
By translation invariance of the ground state,
\(\langle|D(x)|^2\rangle_0\) does not depend on \(x\), so the numerator
of Proposition 1 is \(V\langle|D|^2\rangle_0\), while
\(\langle O^2\rangle_0=\int_V\int_V\langle\varphi(y)\varphi(y')\rangle^{\rm c}_0=V\chi_\varphi\)
with
\(\chi_\varphi=\int d^3z\,\langle\varphi(z)\varphi(0)\rangle^{\rm c}_0\).
The volume cancels:

\textbf{Corollary 2.} For every such \(\varphi\) and every \(V\),
\[\Delta_V\ \le\ \frac{g^2\hbar c}{2}\;\frac{\langle|D_\varphi|^2\rangle_0}{\chi_\varphi}\,,\]
and if \(\Delta_V\to m\) and both ground-state quantities converge, the
same bound holds for \(m\).

Dimensions: \(\delta/\delta A(x)\) carries \(L^{-3}/[A]=L^{-2}\), so
\([\langle|D|^2\rangle]=[\varphi]^2L^2\) and \([\chi]=[\varphi]^2L^3\);
the ratio is an inverse length and \(g^2\hbar c\) times it is an energy.
The plaquette of \href{lattice-gap-upper-bounds.md}{the upper-bound
note} fails precisely because it is a single cell rather than a
zero-momentum sum: its variance is an ultraviolet quantity that vanishes
in the continuum limit while its gradient does not.

\hypertarget{spectral-reading-and-two-checks}{%
\subsection{3. Spectral reading and two
checks}\label{spectral-reading-and-two-checks}}

Inserting a complete set of eigenstates,
\[\langle\Omega,[O,[H,O]]\Omega\rangle=2\sum_n(E_n-E_0)\,\big|\langle n|O|\Omega\rangle\big|^2,
\qquad\langle O^2\rangle_0=\sum_n\big|\langle n|O|\Omega\rangle\big|^2,\]
so Proposition 1 states that \(m\) is at most the mean excitation energy
weighted by the spectral weight of \(O\): the \(f\)-sum rule. Two
consequences.

\emph{Conserved observables.} If \([H,O]=0\) the numerator vanishes and
the bound reads \(\Delta\le0\). With a unique ground state this is
consistent only if \(\langle O^2\rangle_0=0\): a conserved
gauge-invariant observable cannot fluctuate in a nondegenerate vacuum,
since \(O\Omega\) would be a second state of energy \(E_0\). The bound
therefore detects vacuum degeneracy, and is vacuous exactly in that
case.

\emph{Free Maxwell.} Here \(m=0\) and the bound is a positive number for
every flow time; Section 5 shows it tends to zero as the smearing grows,
which is the correct behaviour for a gapless theory and a check that the
bound is not empty.

\hypertarget{what-must-be-smeared-and-the-hypothesis}{%
\subsection{4. What must be smeared, and the
hypothesis}\label{what-must-be-smeared-and-the-hypothesis}}

A local operator at a point has divergent gradient in the continuum, so
\(\varphi\) must be built from a smeared field. The gauge-covariant
smearing is the gradient flow: on the spatial field at fixed time,
\[\partial_s B_i=D_jG_{ji},\qquad B_i|_{s=0}=A_i,\qquad
G_{ij}=\partial_iB_j-\partial_jB_i+[B_i,B_j],\] the three-dimensional,
equal-time counterpart of Lüscher's flow (arXiv:1006.4518v3, equations
(1.1)--(1.2), passage level via the
\href{../docs/Luscher_WilsonFlow_1006.4518v3.md}{local companion}), with
smearing radius \(\sqrt{8t}\) at flow time \(t\) and the lattice form
given by his equation (1.4). Take
\(\varphi=\varphi_t=\tfrac14G^a_{ij}G^a_{ij}\big|_{s=t}-\langle\cdot\rangle_0\).

Lüscher's renormalization statement is for the four-dimensional flow in
the Euclidean functional integral: expectation values of local
gauge-invariant expressions in the flowed field are finite without
further renormalization. The object needed here is the equal-time
Hamiltonian correlator of the spatially flowed field, which his theorem
does not cover. State it as a hypothesis:

\begin{quote}
\textbf{(F1)} For fixed flow time \(t>0\), the equal-time ground-state
quantities \(\chi_{\varphi_t}\) and
\(\langle|D_{\varphi_t}|^2\rangle_0\) of the lattice theory converge, as
\(a\to0\) along the scaling curve, to finite limits, with
\(\chi_{\varphi_t}>0\).
\end{quote}

On any fixed lattice both quantities are finite and positive and
Corollary 2 holds unconditionally, so (F1) is entirely a statement about
the continuum limit. Its first half is the flow's renormalization
property in the Hamiltonian framework; its second half, \(\chi>0\), is
nontriviality: the vacuum is not an eigenstate of the flowed energy
density.

\textbf{Proposition 3 (the finiteness half).} Assume the continuum
theory exists with a unique vacuum and translation invariance, that the
finite-volume gaps converge to \(m\), and (F1) for one flow time \(t\).
Then
\[m\ \le\ \frac{g^2\hbar c}{2}\,\frac{\langle|D_{\varphi_t}|^2\rangle_0}{\chi_{\varphi_t}}\ <\ \infty .\]

This is the \(m<\infty\) clause of the Jaffe--Witten statement, reduced
to the positivity of one susceptibility. The remaining clause, \(m>0\),
is untouched: nothing above excludes \(m=0\).

\hypertarget{the-flow-time-and-where-the-perturbative-scaling-must-fail}{%
\subsection{5. The flow time, and where the perturbative scaling must
fail}\label{the-flow-time-and-where-the-perturbative-scaling-must-fail}}

At flow time \(t\) the only length in the smeared operator is
\(\sqrt{8t}\). In a regime where the flowed correlators are governed by
that length alone, \([\varphi_t]=L^{-4}\) gives
\(\chi\simeq k_1t^{-5/2}\) and \(\langle|D|^2\rangle\simeq k_2t^{-3}\),
hence
\[\frac{\langle|D_{\varphi_t}|^2\rangle}{\chi_{\varphi_t}}\simeq\frac{k_2}{k_1}\,\frac1{\sqrt{8t}},
\qquad m\ \le\ C\,g^2\,\frac{\hbar c}{\sqrt{8t}} .\] The pure numbers
\(k_1,k_2\) are not computed here. Two readings.

\emph{Free theory.} In free Maxwell theory the scaling holds for every
\(t\) because there is no other length, and the bound tends to \(0=m\):
correct, and the flow time is the only resolution scale.

\emph{Interacting theory.} A positive \(m\) cannot be beaten by any
\(t\), so the perturbative scaling must fail once
\(C g^2(t)\hbar c/\sqrt{8t}\) reaches \(m\), that is at
\[\sqrt{8t_*}\ \simeq\ C g^2\,\frac{\hbar c}{m},\] the correlation
length up to the coupling. The bound therefore contains its own
consistency condition: the flow smooths the field only down to the
confinement scale, beyond which the flowed correlators are no longer
controlled by \(\sqrt{8t}\). This is the Hamiltonian counterpart of the
crossover recorded in \href{torus-valley-potential.md}{the valley note},
where the small-volume description fails when the zero-mode energy
reaches \(\hbar c/L\).

\emph{Approach to \(m\).} As \(t\) grows the smeared operator couples
preferentially to states of spatial extent \(\gtrsim\sqrt{8t}\), so the
spectral weights of Section 3 concentrate on the lowest states and the
bound is expected to decrease toward \(m\). Monotonicity is not proved
here; it would make the family of bounds a convergent sequence rather
than a single estimate.

\hypertarget{status}{%
\subsection{6. Status}\label{status}}

\begin{longtable}[]{@{}
  >{\raggedright\arraybackslash}p{(\columnwidth - 2\tabcolsep) * \real{0.5000}}
  >{\raggedright\arraybackslash}p{(\columnwidth - 2\tabcolsep) * \real{0.5000}}@{}}
\toprule\noalign{}
\begin{minipage}[b]{\linewidth}\raggedright
claim
\end{minipage} & \begin{minipage}[b]{\linewidth}\raggedright
status
\end{minipage} \\
\midrule\noalign{}
\endhead
\bottomrule\noalign{}
\endlastfoot
Proposition 1, the bound with explicit constants & proved \\
Corollary 2, intensive zero-momentum form & proved \\
\(f\)-sum rule reading, conserved-observable check & proved \\
finiteness of the lattice bound at fixed \(a\) & proved \\
(F1), continuum limit of the flowed equal-time correlators & hypothesis;
the flow's renormalization in the Hamiltonian framework \\
\(\chi>0\) & hypothesis; nontriviality \\
\(\langle|D|^2\rangle/\chi\simeq c/\sqrt{8t}\) & dimensional,
coefficients not computed \\
\(m>0\) & untouched \\
\end{longtable}

The single-mode bound and the use of smeared operators as glueball
interpolators are standard; the statement being recorded is the
reduction of the Millennium finiteness clause to the positivity of one
flowed susceptibility, with the constants and hypotheses written out.

\hypertarget{consequence-for-state}{%
\subsection{7. Consequence for STATE}\label{consequence-for-state}}

The upper side of the problem is now in the same shape as the lower
side: an exact inequality at fixed cutoff, plus one renormalization
hypothesis for the continuum. The two hypotheses are different in kind.
(F1) concerns flowed observables, which are expected to be renormalized
by Lüscher's four-dimensional result and whose Hamiltonian counterpart
is a bounded, self-contained question; the dressed-vacuum problem of
\href{schur-error-ultraviolet.md}{the Schur note} concerns the state
itself. The next steps this suggests, in order: compute \(k_1,k_2\) in
free Maxwell theory to make Section 5 quantitative and check the free
limit; then state the Hamiltonian flow-renormalization question as its
own target, since both the upper-side (F1) and the lower-side dressing
are instances of it.
\end{document}
